\scanpage{74}
On travaille dans le topos ss $(SS)=\widehat{\Delta}$, avec $k$ anneau comm.

\[
\begin{aligned}
\Phi_*&=\Phi_*k=i_{0*}(k)\\
&=\Phi_{*\mathbb Z}\otimes_{\mathbb Z}k\\
&=\Phi_{*\mathbb Z}\otimes_{\mathbb Z_*}k_* .
\end{aligned}
\]
\deleted{\unsure{J}} on le considère comme faisceau de $k$-modules, (pas comme faisceau d'anneaux) muni d'une section $T$
\[
\left\{\begin{aligned}
T&\in\Gamma(SS,\Phi_*k)\inserted{[=\Gamma(SS,C_*^{10})]}\\
k_*&\hookrightarrow\Phi_*k
\end{aligned}\right.
\quad\text{définissant un \uline{\unsure{régulier}}}.
\]
On a la $k$-Algèbre bigraduée \deleted{\unsure{gr.}} différentielle à p.d.
\[
C_{*k}^{**}\unsure{:=|=}C_{*k}^{**}
=\deleted{\unsure{\Lambda}}\inserted{\Gamma_k^*}\Phi_*\otimes_k\Lambda^*\Phi_*
\]
\margin{(diff. de bidegré $-1,+1$ $\cdots$) \deleted{\unsure{alg.}} alg. \unsure{alternée} pour degré ext,}
qui est une résolution de $k$. De façon précise,
$C_{*k}^{00}=k_*$, et en degré \unsure{total|tot.} donné $n>0$,
on trouve une suite exacte
\[
\begin{gathered}
0-\Gamma^n\Phi^*-\Gamma^{n-1}\Phi^*\otimes\Lambda^1\Phi^*
-\Gamma^{n-2}\Phi^*\otimes\Lambda^2\Phi^*\longrightarrow\\
\cdots\longrightarrow\Phi^*\otimes\Lambda^{n-1}\Phi^*
\longrightarrow\Lambda^n\Phi^*\longrightarrow0\longrightarrow0\cdots
\end{gathered}
\]
dont les termes sont donc flasques.
\margin{Notons $\bar\Phi_*\unsure{\mathrel{\vdash}|:=}\Lambda^1\Phi_*$
et $\bar x$ pour \deleted{\unsure{de}} l'élément \deleted{\unsure{de}} \unsure{de} $\bar\Phi_*$
\unsure{correspondant à} \deleted{\unsure{l'élément}} $x$ de $\Phi_*$.}

On considère $\underset{\Vert\,\bar T}{dT}\in
\Gamma(SS,\deleted{\unsure{\Phi_*}}C_*^{01})\simeq\Gamma(SS,\bar\Phi_*)$
et l'homomorphisme
\[
\bar T_{\wedge}:C_*^{**}\longrightarrow C_k^{**},
\]
on voit qu'il est de bidegré $(0,1)$, \inserted{de carré nul, $\operatorname{Im}(\bar T_\wedge)=\operatorname{Ker}(\bar T_\wedge)$}
\inserted{anti}commutant à $d$
\[
\bigl(d(\bar T\cdot x)=\underset{\Vert\,0}{d\bar T}\cdot x-\bar T\,dx=-\bar T\,dx\bigr),
\]
est injectif sur $C^{*0}$. On pose
\[
\begin{aligned}
DR_{*k}^{pq}=DR_*^{pq}
&=\operatorname{Ker}(C_*^{p,q+1}\xrightarrow{\bar T_\wedge}C_*^{p,q+2})\\
&=\operatorname{Ker}(\Gamma^p\Phi_*\otimes\Lambda^{q+1}\Phi_*-\Gamma^p\Phi_*\otimes\Lambda^{q+2}\Phi_*)\\
&\simeq\deleted{\unsure{\operatorname{Ker}(}}\Gamma^p\Phi_*\otimes
\underbrace{\operatorname{Ker}(\Lambda^{q+1}\Phi_*-\Lambda^{q+2}\Phi_*)}_{\Lambda^q\Psi_*},
\qquad\Psi_*=\Phi_*/k_*\\
&\simeq\Gamma^p\Phi_*\otimes\Lambda^q\Psi_*,
\end{aligned}
\]
$DR_*^{**}=(DR_*^{pq})_{(p,q)\in\mathbb Z\times\mathbb Z}$, avec la bigraduation et la différentielle \inserted{(au signe près)} induites par $C_*^{**}$,
\deleted{\unsure{sauf que}} (mais avec décalage d'une unité pour le degré ext), et une structure multiplicative \deleted{qui} et à p.d. qui sont compatibles
\margin{73}
